\textbf{HP model and designability.} The HP model of Lau and Dill~\cite{lau1989lattice} reduces proteins to H/P chains on a lattice with an energy that rewards H--H contacts; finding the minimum-energy conformation is NP-complete on the 3D cubic lattice~\cite{berger1998protein} and on the 2D square lattice~\cite{crescenzi1998complexity}. Exact enumeration of HP sequences and their designability (for 2D square and diamond lattices) has been studied before~\cite{narasimhan2012hp}, so the counts we compute are not new; what we add is a learned conditional designer and a comparison with search at equal budget. Li et al.~\cite{li1996emergence} studied how many sequences fold to each structure in a simple lattice model and found highly uneven designability; our setting is the inverse of structure prediction on the same model, with exhaustive enumeration playing the role of the oracle.

\textbf{Learned HP folding.} Recent work addresses the \emph{forward} problem of finding the optimal conformation of a given HP sequence with a deep Q-network~\cite{yang2022applying} or with dilated recurrent networks combined with variational annealing~\cite{khandoker2025lattice}. Neither designs sequences for a target. The variational-annealing work~\cite{khandoker2025lattice} combines autoregressive neural models with an annealing schedule and reports competitive folding on the lattice, which motivates our use of simulated annealing as the search baseline.

\textbf{Neural inverse folding.} ProteinMPNN~\cite{dauparas2022robust} and ESM-IF1~\cite{hsu2022learning} generate sequences autoregressively conditioned on backbone structure, and are validated by refolding with predictors such as AlphaFold~\cite{jumper2021highly}. Our designer follows the same pattern at toy scale (a conditional autoregressive model on top of a transformer~\cite{vaswani2017attention}), while replacing the refolding oracle with an exact one.

\textbf{Search-based design.} Simulated annealing~\cite{kirkpatrick1983optimization} is the standard stochastic-search baseline; we reimplement it and show that its performance in this problem depends heavily on the shape of the objective (Section~\ref{sec:results}).
